\chapter{Ретроспективная коррекция меню моделей спектрального рождения}
\label{ch:v6-spectral-transition-candidate-menu-retrospective-correction-gate}

\section{Причина коррекции}

Предыдущий гейт выбрал динамическое поле конечного оператора Дирака как
следующий исследовательский кандидат. Повторный просмотр Тома V показывает,
что это решение не было новым: операторнозначный дефект, его локализованная
фермионная мода, проекторное сокращение кратности и возможные источники
действия уже были последовательно проверены.

Настоящий гейт не стирает положительную кинематику этих результатов, но
отменяет неверное чтение динамического $D_F(x)$ как ещё не исследованной
архитектуры.

\section{Уже построенный динамический массовый дефект}

Гейт Тома V
\path{version5_self_generated_transition_defect_gate}
предъявил нелинейный функционал с kink-профилем, движущимся продолжением и
нормируемой дираковской нулевой модой. Тем самым сама возможность
динамического операторнозначного массового фона уже доказана.

Одновременно были получены два стоп-результата:
\begin{enumerate}
 \item скалярное действие на $M_{20\times15}$ повторяет нулевую моду по
 всем $300$ внутренним направлениям;
 \item двойная яма, её масштаб и отношение Юкавы к жёсткости являются
 новыми входными данными.
\end{enumerate}

\section{Уже проверенное проекторное сокращение}

Следующий гейт
\path{version5_holonomy_projector_defect_multiplicity_gate}
получил условную лестницу
\begin{equation}
 300\longrightarrow45\longrightarrow15\longrightarrow2\longrightarrow1.
\end{equation}
Голономный проектор $P_0$ выбирает семейную линию, лептонный блок --- слабый
дублет, а хиггс-зависимый проектор $P_\nu(H)$ --- нейтринное направление.
Но полный факторный бимодуль допускает только скалярные эндоморфизмы;
$P_\nu(H)$ сингулярен при $H=0$, а амплитуда и единый нелинейный
функционал не выведены.

Следовательно, выражение типа
\begin{equation}
 q(x)P_0\otimes P_\nu(H(x))
 \label{eq:v6-retrospective-operator-defect}
\end{equation}
уже было точной формой открытого операторнозначного дефекта.

\section{Почему спектральное действие не закрывает пробел}

Три независимых результата Тома V закрывают наивное продолжение.

\begin{enumerate}
 \item Гейт
 \path{version5_projector_superconnection_common_scale_gate}
 показал независимость $f_2,f_0,\Lambda,\ell$ и нормировки нечётного поля.
 Один product-Dirac не фиксирует отношение кинетики и потенциала.
 \item Гейт
 \path{spectral_pairing_stiffness_gate}
 показал, что разные допустимые спектральные функции дают разные фазы, а
 вещественное чётное действие не выбирает ориентацию.
 \item Гейт
 \path{version5_h15_fermionic_spectral_weinberg_measure_gate}
 показал независимость чётного и нечётного моментов фермионной спектральной
 функции: кинетическая нормировка не фиксирует амплитуду паринга.
\end{enumerate}

Кроме того,
\path{version5_fermionic_determinant_induced_skyrme_gate}
закрыл попытку получить общий размер через детерминант: коэффициенты
зависят от $M/\Lambda$, юкавских матриц и схемы, а четвёртый порядок не
сводится к единственному положительному члену.

\begin{theorem}[Ретроспективный статус динамического $D_F$]
Динамический конечный оператор Дирака не является новым кандидатом версии
VI. Его локальный дефектный профиль, нулевая мода, проекторное физическое
чтение и спектральные источники коэффициентов уже исследованы в Томе V.
Положительная кинематика сохранена, но общий функционал, масштаб, малая
кратность и устойчивый физический endpoint не выведены. Поэтому прежний
выбор $D_F(x)$ как следующей архитектуры отменяется.
\end{theorem}

\section{Что действительно осталось непостроенным}

Том V прямо зафиксировал, что точный нелинейный дискретный автомат не
построен. Квантовое блуждание дало локальный унитарный перенос и дираковский
непрерывный предел, а непрерывный нелинейный функционал дал kink. Но один
микроскопический закон обновления, который одновременно:
\begin{enumerate}
 \item сохраняет норму и локальность;
 \item имеет ненулевой минимальный размер;
 \item сам порождает два вакуума и дефект;
 \item сокращает внутреннюю кратность без ручного проектора;
 \item выводит непрерывный оператор и его коэффициенты,
\end{enumerate}
ещё не предъявлен.

Это не доказательство существования такого автомата. Это точная граница
между уже закрытым гладким операторнозначным дефектом и действительно
непроверенной дискретной родительской динамикой.

\begin{nogocriterion}
Нельзя переименовывать произвольное поле $D_F(x)$ или заранее заданный kink
в дискретный родитель. Успех требует явного конечного локального правила
обновления, из которого следуют и непрерывный предел, и нелинейность, и
масштаб дефекта. Если потенциал или проектор задаются отдельно, происходит
возврат к уже закрытой модели Тома V.
\end{nogocriterion}

\section{Исправленная архитектурная развилка}

\begin{protocol}
\begin{enumerate}
 \item динамический массовый оператор: \textbf{уже проверен};
 \item kink и нулевая мода: \textbf{получены в proof-of-concept};
 \item кратность полного носителя: \textbf{$300$};
 \item условная проекторная лестница до ранга один:
 \textbf{получена};
 \item единый операторнозначный функционал: \textbf{не получен};
 \item коэффициенты спектрального действия: \textbf{не фиксированы};
 \item фермионный детерминант как общий масштаб: \textbf{закрыт};
 \item прежний выбор $D_F(x)$: \textbf{отменён};
 \item полностью допущенная новая модель: \textbf{отсутствует};
 \item действительно непостроенный объект:
 \textbf{точный нелинейный дискретный родитель}.
\end{enumerate}
\end{protocol}

Следующий гейт
\path{version6_spectral_transition_discrete_nonlinear_parent_reopening_gate}
должен проверить минимальное локальное правило обновления, не вводя
отдельно потенциал, проектор и длину дефекта.

Нулевая гипотеза: унитарность и строгая локальность не позволяют одному
минимальному правилу одновременно дать самофокусирующуюся нелинейность и
физическое сокращение кратности. Стоп-критерий: если для дефекта требуется
независимый нелинейный параметр или внешний проектор, маршрут закрывается
как повтор Тома V.

Машинный сертификат сохранён в
\path{s2t_v6_spectral_transition_candidate_menu_retrospective_correction_gate_results.json}.